\documentclass[11pt,a4paper]{article}
\usepackage[utf8]{inputenc}
\usepackage[T1]{fontenc}
\usepackage{lmodern}
\usepackage{geometry}
\usepackage{amsmath,amssymb,mathtools}
\usepackage{booktabs,longtable,array}
\usepackage{enumitem}
\usepackage{hyperref}
\usepackage{xcolor}
\usepackage{listings}
\usepackage{microtype}
\geometry{margin=25mm}
\hypersetup{colorlinks=true,linkcolor=blue!45!black,urlcolor=blue!45!black}
\lstset{
  language=Python,
  basicstyle=\ttfamily\footnotesize,
  keywordstyle=\color{blue!50!black},
  commentstyle=\color{green!35!black},
  stringstyle=\color{red!45!black},
  breaklines=true,
  frame=single,
  columns=fullflexible,
  showstringspaces=false
}
\DeclareRobustCommand{\code}[1]{\texttt{\detokenize{#1}}}
\newcommand{\shape}[1]{\mathbb{R}^{#1}}
\newcommand{\analogi}[1]{\par\smallskip\noindent\textbf{Analogi.} #1\par\smallskip}
\newcommand{\perhatikan}[1]{\par\smallskip\noindent\textbf{Perhatikan.} #1\par\smallskip}

\title{Membaca dan Mengembangkan Small LLM Cerita\\
\large Panduan kode, rumus, dan contoh langkah demi langkah}
\author{Dokumentasi proyek \code{train-llm-from-scratch}}
\date{23 Oktober 2026}

\begin{document}
\maketitle
\begin{abstract}
Dokumen ini menjelaskan Small LLM dalam proyek ini dari teks mentah sampai
kalimat yang dihasilkan. Sasaran pembacanya adalah orang yang sudah mengenal
Python dasar, tetapi baru mulai mempelajari model bahasa. Setiap rumus
langsung dihubungkan dengan bagian kode dan diterangkan memakai contoh kecil.
Implementasi lokal dan notebook Colab dibahas secara terpisah bila cara
evaluasi atau perangkatnya berbeda. Dokumen ini tidak mengandaikan model sudah
mampu menulis cerita yang baik.
\end{abstract}
\tableofcontents
\newpage

\section{Peta besar: apa yang sedang dipelajari?}
\label{sec:peta}

Small LLM di sini adalah \emph{decoder Transformer} yang belajar menebak
\textbf{byte berikutnya}. Ia tidak menerima arti kata secara langsung. Jika
potongan teks berbunyi "Pagi itu", tugasnya adalah memberi probabilitas
untuk byte yang mungkin muncul sesudahnya. Selama training, tebakan itu
dibandingkan dengan byte yang benar; parameter model lalu diubah sedikit.

\begin{center}
\fbox{\parbox{0.88\linewidth}{\centering
Cerita $\longrightarrow$ byte UTF-8 $\longrightarrow$ jendela input dan target
$\longrightarrow$ embedding $\longrightarrow$ blok Transformer
$\longrightarrow$ skor 257 token $\longrightarrow$ loss
$\longrightarrow$ pembaruan bobot}}
\end{center}

\analogi{Bayangkan murid diminta melanjutkan cerita satu potong sangat kecil
setiap kali. Tokenizer menyiapkan kartu-kartu potongan, attention mengizinkan
murid melihat kartu sebelumnya, dan loss adalah nilai kesalahan dari guru.
Optimizer mengubah kebiasaan murid berdasarkan kesalahan itu.}

\subsection{Istilah yang dipakai}
\begin{description}[leftmargin=35mm,style=nextline]
\item[Token] Satu ID yang diprediksi model. Dalam proyek ini satu token
biasanya satu byte; ID 256 berarti akhir cerita.
\item[Parameter/bobot] Angka yang dipelajari, misalnya isi matriks pada
lapisan linear.
\item[Embedding] Vektor angka yang mewakili token atau posisinya.
\item[Logit] Skor mentah untuk setiap kemungkinan token. Logit belum menjadi
probabilitas.
\item[Loss] Ukuran kesalahan prediksi. Lebih kecil pada data validasi berarti
prediksi byte validasi lebih baik, selama datanya sama.
\item[Gradien] Arah perubahan loss terhadap setiap parameter.
\item[Checkpoint] Salinan bobot dan konfigurasi model untuk inferensi.
\end{description}

\subsection{Peta berkas kode}
\begin{longtable}{p{0.32\linewidth}>{\raggedright\arraybackslash}p{0.60\linewidth}}
\toprule
Berkas & Tanggung jawab \\
\midrule
\endhead
\code{small_llm/tokenizer.py} & Mengubah teks ke byte, mengenali akhir cerita,
dan mengubah byte kembali ke teks. \\
\code{small_llm/data.py} & Membaca satu cerita per berkas, membagi train dan
validasi per cerita, membuat pasangan input--target. \\
\code{small_llm/config.py} & Membaca JSON, memeriksa nilai konfigurasi, memilih
MPS atau CPU pada program lokal. \\
\code{small_llm/model.py} & Embedding, causal attention, blok Transformer,
cross-entropy, dan generasi yang menjaga UTF-8 tetap valid. \\
\code{small_llm/sparse.py} & Lima pilihan MLP: dense, pruned, topk, moe, dan set. \\
\code{small_llm/train.py} & Training satu varian, evaluasi berkala, metrik,
dan checkpoint. \\
\code{small_llm/compare.py} & Menjalankan beberapa varian dan membuat tabel
perbandingan. \\
\code{small_llm/generate.py} & Memuat checkpoint dan melanjutkan prompt. \\
\code{small_llm/__init__.py} & Menandai \code{small_llm} sebagai paket Python;
tidak menjalankan training. \\
\code{tests/test_small_llm.py} & Memeriksa sifat penting kode dengan model
dan data kecil. \\
Notebook Colab & Versi mandiri dalam
\code{small_llm_frog_}\allowbreak\code{storytelling_colab.ipynb};
seluruh kode ditulis di notebook, dengan dataset Hugging Face. \\
\bottomrule
\end{longtable}

\section{Langkah 1: data menjadi token}
\label{sec:data}
\subsection{Tokenizer byte: \code{tokenizer.py}}

UTF-8 mengubah karakter menjadi satu atau beberapa byte. Fungsi
\code{encode(text)} menjalankan \code{text.encode("utf-8")} lalu mengubahnya
menjadi daftar bilangan. Kosakatanya tetap:
\[
\mathcal{V}=\{0,1,\ldots,255,256\},\qquad |\mathcal{V}|=257.
\]
ID 0--255 adalah nilai byte. ID 256 adalah \emph{END\_OF\_STORY}, penanda
akhir dokumen yang ditambahkan saat data dibaca. Contoh sederhana:
\[
\text{"A"}\mapsto[65], \qquad
\text{"AB"}\mapsto[65,66], \qquad
\text{"A" + akhir cerita}\mapsto[65,256].
\]
Huruf beraksen atau simbol tertentu bisa memerlukan lebih dari satu byte.
Karena itu panjang konteks 256 berarti \textbf{256 byte/token}, bukan
256 kata. Fungsi \code{decode(ids)} berhenti pada ID 256. Jika daftar byte
yang diberikan rusak, fungsi itu mengganti bagian UTF-8 yang tidak lengkap
agar teks tetap dapat ditampilkan.

\analogi{Token seperti potongan puzzle. Satu huruf biasa dapat menjadi satu
potong, sementara karakter lain perlu beberapa potong. Kartu 256 adalah
tanda "cerita selesai".}

\perhatikan{Tokenizer byte tidak mempunyai masalah kata "tidak dikenal",
karena setiap teks UTF-8 dapat dipecah menjadi byte. Namun urutan token
menjadi lebih panjang dibanding tokenizer subkata.}

\subsection{Pemisahan cerita: \code{data.py}}

\code{load_stories} membaca semua berkas \code{*.txt}, mengurutkan namanya,
lalu mengacak dengan \code{seed}. Sebelum token digabung, berkas dipisah
menjadi train dan validasi. Untuk $N$ berkas dan fraksi validasi $r$, kode
menggunakan
\[
N_{\mathrm{valid}}=
\max\!\left(1,\min\!\left(N-1,\operatorname{round}(Nr)\right)\right),
\qquad
N_{\mathrm{train}}=N-N_{\mathrm{valid}}.
\]
Sebagai contoh, $N=60$ dan $r=0{,}25$ menghasilkan 15 cerita validasi dan
45 cerita train. Pemisahan per berkas mencegah awal dan akhir cerita yang
sama masuk ke kedua kelompok. Setelah dibagi, setiap kelompok digabung
menjadi satu aliran token, dengan ID 256 di antara cerita. Karena itu satu
jendela dapat melintasi batas dua cerita; batasnya terlihat oleh model melalui
ID 256.

\analogi{Jangan membagi halaman satu buku ke dua ruang ujian. Simpan seluruh
buku di ruang train atau seluruhnya di ruang validasi. Sesudah pembagian,
halaman-halaman dalam tiap ruang baru disusun menjadi tumpukan.}

\subsection{Pasangan input--target: \code{sample_batch}}

Misalkan aliran token adalah $C=(c_0,c_1,\ldots,c_{N-1})$. Untuk panjang
konteks $L$, sebuah titik awal $s$ menghasilkan
\[
X=(c_s,c_{s+1},\ldots,c_{s+L-1}),\qquad
Y=(c_{s+1},c_{s+2},\ldots,c_{s+L}).
\]
Target $Y$ bergeser satu posisi dari input $X$. Jika alirannya
$[10,20,30,40,50]$ dan $L=3$, salah satu pasangan adalah
$X=[10,20,30]$ dan $Y=[20,30,40]$. Model belajar bahwa setelah 10
muncul 20, setelah 20 muncul 30, dan seterusnya. \code{batch_size=B}
berarti ada $B$ jendela seperti ini dalam satu microbatch. Bentuk tensor
\[
X,Y\in\mathbb{Z}^{B\times L}.
\]
Kode memilih $s$ secara acak dari aliran yang sesuai dan memindahkan batch
ke \code{device}.

\subsection{Konfigurasi: \code{config.py} dan tiga JSON}

\code{ModelConfig} menyimpan $L$ (\code{context_length}), $d$
(\code{n_embed}), jumlah head $H$, dan jumlah blok $K$.
\code{TrainConfig} menyimpan data, output, seed, batch, langkah, dan
pengaturan optimizer. \code{SparseConfig} menyimpan pilihan untuk tiap
varian MLP. \code{load_config} membaca JSON dan membangun tiga kelompok
pengaturan itu. Nilai yang tidak sah, misalnya $d$ tidak habis dibagi $H$,
ditolak sebelum training.

\begin{center}
\begin{tabular}{p{0.35\linewidth}p{0.45\linewidth}}
\toprule
Konfigurasi & Kegunaan \\
\midrule
\code{configs/smoke.json} & Dua langkah per varian; memeriksa alur, bukan
kualitas cerita. \\
\code{configs/small_m1.json} & Model $d=192$, $H=4$, $K=4$, $L=256$ pada
12 cerita demo. \\
\code{stories_expanded_m1.json} & 60 cerita, fraksi validasi
$0{,}25$, 600 langkah; output terpisah. \\
\bottomrule
\end{tabular}
\end{center}

Pada program lokal, \code{device="auto"} memilih MPS jika tersedia dan CPU
jika tidak. Notebook Colab memilih CUDA jika GPU aktif, selain itu CPU.
Mengubah ukuran model membuat checkpoint lama tidak cocok dengan arsitektur
baru.

\section{Langkah 2: Transformer membuat prediksi}
\label{sec:model}
\subsection{Embedding token dan posisi}

Untuk token $t_i$ pada posisi $i$, model mengambil dua vektor berdimensi $d$:
\[
\mathbf{x}_i=E_{\mathrm{token}}[t_i]+E_{\mathrm{pos}}[i],
\quad E_{\mathrm{token}}\in\shape{257\times d},
\quad E_{\mathrm{pos}}\in\shape{L\times d}.
\]
Penjumlahan ini dilakukan oleh \code{token_embedding} dan
\code{position_embedding}. Vektor token memberi identitas potongan teks;
vektor posisi memberi tahu di mana ia berada dalam jendela. Bentuk hasilnya
$B\times L\times d$.

\analogi{Kartu bertuliskan "pagi" memberi isi, nomor kursi memberi posisi.
Dua informasi itu diperlukan agar murid dapat membedakan "pagi itu" dari
"itu pagi".}

\subsection{Causal attention: \code{CausalAttention}}

Setiap head memproyeksikan vektor menjadi query, key, dan value:
\[
Q=XW_Q,\qquad K=XW_K,\qquad V=XW_V,\qquad d_h=d/H.
\]
Untuk posisi $i$ yang ingin memprediksi token berikutnya, skor melihat posisi
$j$ adalah
\[
a_{ij}=\frac{\mathbf{q}_i^\top\mathbf{k}_j}{\sqrt{d_h}}+M_{ij},
\qquad
M_{ij}=
\begin{cases}
0,&j\le i,\\
-\infty,&j>i.
\end{cases}
\]
Softmax mengubah skor menjadi bobot perhatian
$\alpha_{ij}=\exp(a_{ij})/\sum_{k\le i}\exp(a_{ik})$. Keluaran head pada
posisi $i$ adalah $\mathbf{o}_i=\sum_{j\le i}\alpha_{ij}\mathbf{v}_j$.
Keluaran semua head digabung dan diproyeksikan lagi.

\analogi{Query adalah pertanyaan yang diajukan kata saat ini; key adalah
label pada catatan lama; value adalah isi catatan. Aturan \emph{causal}
menutup halaman-halaman masa depan. Tanpanya, model bisa mengintip jawaban
saat latihan.}

\perhatikan{Kode lokal membuat mask segitiga bawah dengan
\code{torch.tril}. Notebook Colab memakai operasi
\code{scaled_dot_product_attention(..., is_causal=True)}; tujuan matematisnya
sama. Untuk $d=192$ dan $H=4$, tiap head memakai $d_h=48$ dimensi.}

\subsection{Residual, LayerNorm, dan MLP}

Satu \code{Block} menjalankan dua operasi berurutan:
\begin{align}
A &= X+\operatorname{Attention}(\operatorname{LN}(X)),\\
X_{\mathrm{next}} &= A+\operatorname{MLP}(\operatorname{LN}(A)).
\end{align}
Tanda $+$ adalah \emph{residual connection}: informasi lama tetap punya
jalur langsung. \code{LayerNorm} menstabilkan skala vektor. MLP dense
mengembangkan lebar dari $d$ ke $4d$, menerapkan ReLU, lalu kembali ke $d$:
\[
\operatorname{MLP}(\mathbf{x})=
W_2\,\operatorname{ReLU}(W_1\mathbf{x}+\mathbf{b}_1)+\mathbf{b}_2,
\quad
W_1\in\shape{4d\times d},\quad
W_2\in\shape{d\times 4d}.
\]
Untuk $d=192$, lebar tersembunyi MLP adalah 768. Blok diulang $K=4$
kali pada konfigurasi M1.

\analogi{Attention mengumpulkan catatan yang relevan dari bagian cerita
sebelumnya. MLP memikirkan ulang catatan itu untuk posisi saat ini.
Residual seperti menyimpan salinan catatan asli agar tidak hilang saat
ditafsirkan.}

\subsection{Logit, probabilitas, dan prediksi berikutnya}

Setelah blok terakhir, \code{final_norm} dan \code{lm_head} menghasilkan
logit $z_{i,v}$ untuk setiap $v\in\mathcal V$:
\[
Z=\operatorname{Linear}(\operatorname{LN}(X_{\mathrm{akhir}})),
\qquad Z\in\shape{B\times L\times257}.
\]
Probabilitas token $v$ pada posisi $i$ diperoleh melalui softmax,
\[
p(v\mid t_{\le i})=
\frac{\exp(z_{i,v})}{\sum_{u=0}^{256}\exp(z_{i,u})}.
\]
Logit besar berarti peluang besar relatif terhadap token lain; logit
sendiri bukan persentase.

\section{Langkah 3: lima pilihan MLP dalam \code{sparse.py}}
\label{sec:sparse}

Semua varian memakai tokenizer, attention, data, dan prosedur training yang
sama. Setelah Transformer dense dibuat, \code{convert_model} mengganti
MLP pada tiap blok.

\subsection{Dense: acuan}

Dense memakai seluruh $4d$ neuron dan seluruh bobot $W_1,W_2$. Nilai loss,
kecepatan, dan memori dense menjadi acuan untuk empat varian lain.
\analogi{Semua pekerja di ruang kerja mendapat tugas pada setiap token.}

\subsection{Pruned: mask bobot tetap}

Kelas \code{PrunedMLP} mencari bobot dengan nilai mutlak terkecil saat
model dibentuk. Untuk target sparsity $s$, sekitar fraksi $s$ elemen pada
tiap matriks diberi mask nol. Operasi yang dipakai saat \code{forward}:
\[
W_{\mathrm{efektif}}=W\odot M,\qquad
M_{ij}\in\{0,1\}.
\]
Simbol $\odot$ berarti perkalian elemen demi elemen. Posisi dengan $M_{ij}=0$
tidak berpengaruh pada keluaran. Mask disimpan sebagai buffer, bukan parameter
yang dilatih. Bobot matriks tetap dialokasikan penuh; implementasi ini tidak
menjamin RAM atau waktu lebih kecil.

\analogi{Sebuah peta jalan tetap dicetak lengkap, tetapi beberapa ruas diberi
tanda tertutup. Menutup ruas bukan berarti ukuran kertas petanya mengecil.}

\subsection{TopK: beberapa neuron dipakai per token}

\code{TopKMLP} terlebih dahulu menghitung
$\mathbf{h}=\operatorname{ReLU}(W_1\mathbf{x}+\mathbf b_1)$.
Dari $4d$ aktivasi, hanya $k=\max(1,\operatorname{round}(4df))$ yang
terbesar dipilih, dengan $f=\code{active_fraction}$. Keluaran:
\[
\mathbf y=\mathbf b_2+\sum_{j\in\operatorname{TopK}(\mathbf h,k)}
h_j W_{2,:,j}.
\]
Dengan $d=192$ dan $f=0{,}25$, $4d=768$ dan $k=192$ neuron dipakai
untuk setiap token. Lapisan pertama tetap menghitung seluruh 768 aktivasi.
Pemilihan \code{topk} dan pengumpulan kolom dapat memiliki biaya sendiri;
angka neuron aktif tidak otomatis berarti akselerasi empat kali.

\analogi{Semua pekerja mengajukan ide awal, lalu hanya ide dengan skor
tertinggi dibawa ke rapat akhir. Biaya mengumpulkan ide awal masih ada.}

\subsection{MoE: memilih beberapa expert}

\code{MoEMLP} membagi lebar MLP dense ke $E$ expert. Router atau
\emph{gate} memberi logit $g_e$ untuk expert $e$ dan memilih himpunan
$S=\operatorname{TopK}(\mathbf g,k)$:
\[
\alpha_e=
\frac{\exp(g_e)}{\sum_{j\in S}\exp(g_j)}
\quad(e\in S),\qquad
\mathbf y=k\sum_{e\in S}\alpha_e\,\operatorname{Expert}_e(\mathbf x).
\]
Faktor $k$ berasal dari implementasi ini untuk menjaga skala ketika
beberapa expert dijumlahkan. Nilai \code{top_k=2} dan
\code{n_experts=4} berarti dua dari empat expert dipakai per token.
Expert mula-mula menerima potongan bobot MLP dense, sementara router
dilatih dari awal.

Tanpa penyeimbang, router dapat terlalu sering memilih expert yang sama.
Kode menambahkan loss bantu:
\[
\mathcal L_{\mathrm{aux}}
=E\sum_{e=1}^{E} I_e L_e,
\qquad
I_e=\operatorname{mean}(\operatorname{softmax}(\mathbf g)_e),
\quad
L_e=\operatorname{mean}(\mathbf 1[e\in S]).
\]
Loss total saat training ialah
$\mathcal L_{\mathrm{token}}+\lambda\mathcal L_{\mathrm{aux}}$.
Nilai $\lambda$ diatur oleh \code{aux_coef}.
Metrik \code{train_loss} dan \code{valid_loss} hanya melaporkan
loss token agar prediksi dapat dibandingkan.

\analogi{Router adalah petugas penerima naskah. Ia mengirim satu potongan
cerita ke dua editor dari empat editor yang tersedia. Loss bantu mencegah
satu editor terus-menerus kebanjiran pekerjaan.}

\subsection{SET: mask berubah selama training}

\code{SETMLP} membuat mask acak dengan kerapatan awal
$\rho=\code{density}$, sehingga sekitar $\rho$ fraksi koneksi aktif.
Selama training, bobot bermask nol dipaksa tetap nol. Setiap
\code{rewire_every} langkah, \code{rewire_set} melakukan:
\begin{enumerate}
\item Cari koneksi aktif dengan $|w|$ terkecil dan matikan sejumlah
$r=\operatorname{round}(\code{rewire_fraction}\times N_{\mathrm{aktif}})$.
\item Pilih $r$ posisi tidak aktif secara acak dan hidupkan.
\item Nolkan bobot pada posisi yang berubah dan state optimizer yang
bentuknya sama dengan matriks bobot.
\end{enumerate}
Jumlah koneksi aktif tetap, tetapi lokasinya berubah. Menolkan state
optimizer penting karena momentum lama tidak boleh terbawa ke sambungan
baru.

\analogi{Jumlah kursi dalam tim tetap. Secara berkala kursi dengan
kontribusi terendah ditutup dan kursi baru dibuka di tempat lain. Catatan
kerja lama pada kursi yang berpindah dihapus.}

\perhatikan{\code{active_connections} menghitung elemen mask aktif pada
\code{pruned} dan \code{set}. Ukuran itu tidak sama dengan jumlah parameter
fisik atau banyaknya operasi GPU untuk \code{topk} dan \code{moe}.}

\section{Langkah 4: kesalahan mengubah bobot}
\label{sec:training}
\subsection{Cross-entropy: \code{model.py}}

Untuk target benar $y_{b,i}$ pada batch $b$, posisi $i$, model memperoleh
probabilitas $p(y_{b,i}\mid X_{b,\le i})$. Loss rata-rata:
\[
\mathcal L_{\mathrm{token}}=
-\frac{1}{BL}\sum_{b=1}^{B}\sum_{i=1}^{L}
\log p(y_{b,i}\mid X_{b,\le i}).
\]
Contoh satu tebakan: jika token benar diberi peluang $0{,}5$, kesalahannya
$-\log(0{,}5)\approx0{,}693$. Jika peluangnya hanya $0{,}01$,
kesalahannya $-\log(0{,}01)\approx4{,}605$. Kode memakai
\code{F.cross_entropy} pada logit yang diratakan ke bentuk
$(BL)\times257$ dan target sebanyak $BL$.

Jika seluruh 257 token diberi peluang sama, loss-nya
$\log(257)\approx5{,}55$. Itu acuan kasar untuk tebakan seragam,
\textbf{bukan} target mutu cerita. Loss rendah juga bisa terjadi karena
model menguasai ejaan dan tanda baca tanpa memahami alur.

\analogi{Guru menghukum tebakan yang yakin tetapi salah lebih besar
daripada tebakan yang masih memberi peluang wajar pada jawaban benar.}

\subsection{Akumulasi gradien dan optimizer: \code{train.py}}

Satu langkah optimizer terdiri atas $G=\code{grad_accum}$ microbatch.
Kode membagi loss tiap microbatch dengan $G$ sebelum \code{backward}.
Maka gradien yang terkumpul setara dengan rata-ratanya:
\[
\mathbf g
=\frac{1}{G}\sum_{j=1}^{G}
\nabla_\theta\mathcal L_j.
\]
Jumlah posisi token yang diproses per langkah adalah
$BGL$. Pada konfigurasi M1: $B=2$, $G=4$, $L=256$, sehingga
$BGL=2048$ posisi per langkah. Dalam 300 langkah, model memproses
$614\,400$ posisi, tetapi banyak jendela diambil berulang dari korpus
yang jauh lebih kecil. Jumlah posisi diproses tidak sama dengan jumlah
byte unik.

Setelah \code{backward}, kode memotong norma gradien pada
\code{grad_clip=c}:
\[
\mathbf g_{\mathrm{clip}}=
\mathbf g\min\!\left(1,\frac{c}{\|\mathbf g\|_2}\right).
\]
Kemudian \code{AdamW} memakai riwayat gradien untuk memperbarui
parameter. Bentuk ringkas satu langkahnya:
\begin{align}
\mathbf m_t&=\beta_1\mathbf m_{t-1}+(1-\beta_1)\mathbf g_t,\\
\mathbf v_t&=\beta_2\mathbf v_{t-1}+(1-\beta_2)\mathbf g_t^2,\\
\theta_{t+1}
&\approx (1-\eta\lambda_{\mathrm{wd}})\theta_t
-\eta\frac{\widehat{\mathbf m}_t}
{\sqrt{\widehat{\mathbf v}_t}+\epsilon}.
\end{align}
Di sini $\eta$ adalah \code{learning_rate}; $\lambda_{\mathrm{wd}}$
adalah \code{weight_decay}. Implementasi rinci, termasuk koreksi bias,
ditangani PyTorch.

\analogi{Gradien adalah petunjuk arah turun bukit. Akumulasi meminta
pendapat beberapa kelompok kecil sebelum melangkah. Gradient clipping
mencegah satu pendapat ekstrem membuat langkah terlalu jauh. AdamW
mengingat arah yang sering membantu.}

\subsection{Urutan satu langkah dalam kode}

\begin{enumerate}
\item \code{model.train()} mengaktifkan mode training dan
\code{optimizer.zero_grad()} mengosongkan gradien lama.
\item Ambil $G$ batch. Untuk setiap batch: hitung logit, loss token,
tambahkan loss bantu MoE bila perlu, bagi dengan $G$, lalu
\code{backward()}.
\item Periksa bahwa loss berupa angka hingga; potong gradien; panggil
\code{optimizer.step()}.
\item Jika varian SET, paksa bobot tidak aktif ke nol dan jalankan
rewiring pada intervalnya.
\item Tulis \code{metrics.jsonl}: catat waktu langkah dan, pada interval
tertentu, evaluasi pada data validasi.
\end{enumerate}

\perhatikan{Pada MPS, \code{_sync} menunggu operasi perangkat selesai
sebelum dan sesudah pengukuran waktu. Tanpa sinkronisasi, waktu dapat
terlihat lebih kecil karena tugas GPU masih berjalan.}

\section{Langkah 5: evaluasi dan berkas hasil}
\label{sec:evaluasi}

\subsection{Evaluasi lokal: \code{evaluate} dalam \code{train.py}}

Model masuk ke mode \code{eval} dan berjalan tanpa gradien. Fungsi ini
mengambil \code{eval_batches} jendela acak dari cerita validasi,
kemudian menghitung rata-rata cross-entropy:
\[
\mathcal L_{\mathrm{valid, lokal}}
=\frac{1}{J}\sum_{j=1}^{J}\mathcal L_j,
\qquad J=\code{eval_batches}.
\]
Generator titik awal selalu dibuat lagi dengan \code{seed + 2}. Jadi,
pada setiap evaluasi, jendela validasi yang dipakai sama. Perubahan loss
lebih mudah ditafsirkan karena pertanyaannya tidak berganti-ganti.
Namun yang diukur tetap sampel jendela, bukan semua byte validasi.

\analogi{Murid mengerjakan set soal ujian yang sama setiap beberapa jam.
Kemajuannya terlihat lebih jelas daripada bila soalnya selalu diacak,
tetapi hasil itu belum mewakili semua kemungkinan soal.}

Baris langkah 0 menyimpan loss validasi sebelum pembaruan bobot.
\code{eval_every} menentukan interval evaluasi berikutnya; langkah terakhir
selalu dievaluasi. \code{log_every} menentukan kapan baris training dicatat.
Pada baris tanpa evaluasi, \code{valid_loss} bernilai \code{null}, bukan nol.

\subsection{Cara membaca hasil}

\code{metrics.jsonl} berisi satu objek JSON per baris, misalnya:
\begin{lstlisting}[language={}]
{"step": 100, "train_loss": 2.20, "valid_loss": 2.30,
 "step_seconds": 0.06}
\end{lstlisting}
\code{summary.json} meringkas nilai terakhir, jumlah parameter, pembagian
berkas, jumlah token, waktu, dan memori. \code{model.pt} berisi bobot,
konfigurasi arsitektur, konfigurasi sparsity, dan nama varian. Checkpoint
lokal tidak menyimpan state AdamW atau posisi generator acak; menjalankan
training lagi berarti memulai percobaan baru.

\begin{center}
\begin{tabular}{p{0.32\linewidth}p{0.59\linewidth}}
\toprule
Metrik & Cara membaca \\
\midrule
\code{train_loss} & Loss token pada microbatch langkah terakhir;
berfluktuasi karena batch berubah. \\
\code{valid_loss} & Loss pada jendela validasi yang dipilih dengan seed
tetap; bandingkan dalam eksperimen yang datanya sama. \\
\code{initial_valid_loss} & Titik awal sebelum training. \\
\code{train_tokens_per_second} & $BGL\times\code{steps}$ dibagi waktu
langkah training; evaluasi dan generasi tidak dihitung. \\
\code{peak_process_rss_mb} & Puncak memori proses Python. Ini bukan ukuran
lengkap semua memori GPU atau biaya produksi. \\
\bottomrule
\end{tabular}
\end{center}

Jika \code{train_loss} turun tetapi \code{valid_loss} naik terus, model
mungkin terlalu menyesuaikan diri pada data train. Jika keduanya turun,
model belajar memprediksi token pada distribusi yang diuji. Kualitas
cerita masih harus diperiksa dengan membaca contoh baru.

Perplexity per token dapat dihitung dari cross-entropy:
\[
\operatorname{PPL}=\exp(\mathcal L).
\]
Karena token di sini adalah byte, ini adalah perplexity per byte.
Jangan bandingkan langsung dengan angka dari tokenizer kata atau subkata.
Program lokal tidak menyimpan perplexity sebagai kolom terpisah.

\subsection{Perbandingan: \code{compare.py}}

\code{compare.py} menjalankan dense dan empat varian sparse dengan
konfigurasi yang sama. Secara bawaan tiap varian memakai proses Python
terpisah, agar puncak RSS satu varian tidak terbawa ke pengukuran varian
lain. Setelah itu program mengumpulkan \code{summary.json} menjadi
\code{comparison.csv}. Opsi \code{--variants} memilih sebagian varian;
\code{--overwrite} diperlukan jika folder output varian sudah ada dan
pengguna memang ingin menggantinya. Dengan \code{--in-process}, varian
dijalankan dalam satu proses dan angka puncak RSS per varian dikosongkan.

\perhatikan{Perbandingan adil memerlukan data, pembagian train/validasi,
seed, batch token, jumlah langkah, optimizer, dan perangkat yang sama.
Jumlah parameter serta operasi tiap varian tetap berbeda. Loss dari
\code{stories_demo} tidak langsung sebanding dengan loss dari
\code{stories_expanded} atau dataset Frog karena data ujinya berubah.}

\section{Langkah 6: model menulis lanjutan}
\label{sec:generasi}

\code{generate.py} memuat \code{model.pt} dengan
\code{torch.load(..., weights_only=True)}, membangun kembali Transformer
dan jenis MLP dari metadata, lalu memuat \code{state_dict}. Prompt diubah
menjadi byte, kemudian \code{StoryTransformer.generate} mengulang:
\begin{enumerate}
\item Ambil paling banyak $L$ token terakhir sebagai konteks.
\item Hitung logit pada posisi terakhir.
\item Bagi logit dengan suhu $T>0$, ubah menjadi probabilitas,
pilih satu token, lalu tambahkan ke konteks.
\item Berhenti jika ID 256 terpilih atau batas token baru tercapai.
\end{enumerate}
Distribusi sampling pada suhu $T$ adalah
\[
p_T(v)=
\frac{\exp(z_v/T)}
{\sum_{u\in A}\exp(z_u/T)},\qquad v\in A,
\]
dengan $A$ himpunan token yang diizinkan. Suhu lebih rendah membuat
pilihan lebih terkonsentrasi pada logit tinggi; suhu lebih tinggi
membuat hasil lebih bervariasi. Mengubah suhu tidak mengajari model
fakta atau alur baru.

\subsection{Mengapa keluaran tetap UTF-8 tetapi bisa aneh?}

Karena satu karakter Unicode dapat memakai beberapa byte, sampling byte
secara bebas bisa menghasilkan urutan yang rusak. Implementasi lokal
melacak apakah model sedang berada di tengah karakter UTF-8. Jika ya,
hanya byte lanjutan yang sah diizinkan. Di akhir batas token, model
tidak boleh memulai karakter yang tidak akan sempat diselesaikan.

\analogi{Untuk menulis satu huruf yang terdiri dari tiga potong puzzle,
model harus menyelesaikan potong kedua dan ketiga sebelum memulai huruf
baru. Aturan ini menjaga bentuk tulisannya, bukan isi ceritanya.}

Generasi menampilkan kembali prompt diikuti kelanjutannya. Batas
\code{--max-new-tokens} menghitung byte baru, bukan kata atau
kalimat. Model 300 langkah pada 12 cerita dapat membuat suku kata yang
menyerupai bahasa Indonesia tetapi belum menghasilkan cerita runtut.

\section{Tes: apa yang dibuktikan \code{test_small_llm.py}?}
\label{sec:tes}

\begin{enumerate}
\item \textbf{Tokenizer dan data:} teks UTF-8 dapat dikodekan dan dibaca
kembali; berkas train dan validasi tidak tumpang tindih; target adalah
input yang bergeser satu token.
\item \textbf{Causal attention:} mengubah token masa depan tidak mengubah
logit posisi sebelumnya.
\item \textbf{TopK dan MoE:} dalam pengaturan yang memakai semua
neuron/expert, keluaran dapat disamakan dengan MLP dense sesuai
konstruksi tes.
\item \textbf{SET:} rewiring menjaga jumlah koneksi aktif dan menghapus
state optimizer pada posisi yang berubah.
\item \textbf{Checkpoint:} arsitektur sparse dapat dibangun kembali
dan menghasilkan keluaran yang sama sesudah bobot dimuat.
\item \textbf{Training singkat:} setiap varian dapat menjalankan satu
langkah dan menyimpan checkpoint dengan loss hingga.
\item \textbf{Generasi:} byte baru mengikuti urutan UTF-8 yang valid
pada contoh yang sengaja memberi skor tinggi pada byte tidak sah.
\end{enumerate}

\perhatikan{Tes ini membuktikan sifat kode pada contoh kecil. Tes satu
langkah tidak membuktikan bahwa model mampu bercerita. Kualitas bahasa
memerlukan data dan evaluasi yang terpisah.}

\section{Versi mandiri Colab: apa yang sama dan berbeda?}
\label{sec:colab}

Notebook \code{small_llm_frog_storytelling_colab.ipynb} menulis ulang
tokenizer, Transformer, kelima varian MLP, training, checkpoint, dan
generasi di dalam satu berkas. Ia tidak mengimpor
\code{small_llm/*.py} atau cerita lokal. Data dibaca dari
\href{https://huggingface.co/datasets/bigscience-data/roots_id_indonesian_frog_storytelling_corpus}{halaman dataset Frog Storytelling di Hugging Face}.
Dataset tersebut mewajibkan persetujuan akses di situs Hugging Face dan
login akun yang sama di Colab. Notebook memeriksa apakah kolom
\code{text} ada dan menampilkan nama kolom jika tidak ada.

\begin{enumerate}
\item Sel pemasangan dependensi dan login menyiapkan pembacaan Parquet dari
alamat \code{hf://}. Token akses tidak ditulis dalam notebook.
\item Sel data membersihkan spasi, membuang baris kosong dan duplikat
identik, lalu membagi dokumen train/validasi dengan seed tetap. Dokumen
baru digabung menjadi aliran byte setelah pembagian.
\item Sel model dan varian menerapkan rumus pada Bagian
\ref{sec:model}--\ref{sec:sparse}. Attention Colab memakai operasi
PyTorch \code{is_causal=True}.
\item Sel training melatih \code{dense} secara bawaan.
\code{RUN_COMPARISON} perlu diubah menjadi \code{True} untuk melatih
empat varian lain.
\item Sel evaluasi membaca semua byte validasi secara berurutan dalam
potongan konteks, bukan beberapa jendela acak seperti \code{train.py}
lokal. Ia menyimpan loss validasi, perplexity per byte, dan akurasi
prediksi byte.
\item Saat loss validasi membaik, notebook menyimpan
\code{best_model.pt}. Model untuk generasi dimuat dari checkpoint
terbaik ini, bukan otomatis dari langkah terakhir.
\item Generasi notebook membatasi karakter pada karakter yang
muncul dalam dokumen train melalui trie byte UTF-8. Aturan ini lebih
ketat daripada mask UTF-8 umum pada program lokal, tetapi tetap tidak
menjamin alur cerita masuk akal.
\item Sel terakhir mengemas hasil ke ZIP agar dapat diunduh sebelum
runtime Colab berhenti.
\end{enumerate}

\perhatikan{Dataset Frog berasal dari penceritaan atas gambar yang sama.
Dokumen validasi dipisah, tetapi topik dan urutan adegan antar cerita
dapat sangat mirip. Angka validasi karenanya tidak boleh dianggap
bukti kemampuan menulis cerita bertopik bebas. Nilai notebook juga
tidak langsung sebanding dengan nilai program lokal: dataset, cara
evaluasi, jumlah langkah, dan checkpoint yang dipilih berbeda.}

\section{Contoh ukuran tensor dari awal sampai akhir}
\label{sec:ukuran}

Dengan konfigurasi \code{small_m1.json}, ambil $B=2$, $L=256$,
$d=192$, $H=4$, dan $V=257$. Ikuti satu batch:
\begin{center}
\begin{tabular}{p{0.45\linewidth}p{0.35\linewidth}}
\toprule
Tahap & Bentuk tensor \\
\midrule
Input dan target & $2\times256$ \\
Embedding token + posisi & $2\times256\times192$ \\
Query/key/value sebelum dipisah head & masing-masing
$2\times256\times192$ \\
Query/key/value sesudah dipisah head & masing-masing
$2\times4\times256\times48$ \\
Skor attention per head & $2\times4\times256\times256$ \\
Aktivasi tersembunyi MLP dense & $2\times256\times768$ \\
Logit keluaran & $2\times256\times257$ \\
Target untuk cross-entropy & $512$ ID setelah diratakan \\
\bottomrule
\end{tabular}
\end{center}
Tabel ini berguna saat mengubah arsitektur. Jika sebuah operasi gagal
karena ukuran tensor, cari baris pertama yang bentuknya tidak sesuai.

\section{Langkah pengembangan yang dapat diulang}
\label{sec:pengembangan}

\subsection{Mulai dari satu pertanyaan}

Contoh pertanyaan: apakah konteks lebih panjang membantu loss validasi
dan kesinambungan kalimat? Buat satu baseline, simpan konfigurasi,
seed, pembagian data, \code{metrics.jsonl}, dan beberapa prompt yang
tetap. Pada percobaan berikutnya, ubah satu hal, misalnya
\code{context_length}. Jika sekaligus mengubah dataset dan tokenizer,
sulit mengetahui penyebab perubahan hasil.

\begin{center}
\begin{tabular}{p{0.21\linewidth}p{0.29\linewidth}p{0.39\linewidth}}
\toprule
Ide & Perubahan kode/konfigurasi & Pemeriksaan yang diperlukan \\
\midrule
Tambah data & Ganti \code{data_dir}; bagi per cerita &
Periksa duplikat dan jumlah cerita validasi; jangan bandingkan
loss mentah dengan dataset lama. \\
Konteks lebih panjang & Naikkan \code{context_length} &
Periksa RAM, bentuk mask/posisi, loss, dan kelanjutan cerita. \\
Tokenizer subkata & Ganti \code{encode/decode}, ukuran kosakata,
embedding, dan head keluaran & Latih ulang dari awal; checkpoint
tokenizer byte tidak kompatibel. \\
MLP sparse & Ubah satu opsi sparsity/topk/moe/set &
Catat loss, jumlah parameter, koneksi aktif, waktu, dan memori. \\
Berhenti lebih dini & Simpan checkpoint terbaik menurut
\code{valid_loss} & Pastikan validasi tidak dipakai untuk melatih
bobot; sediakan set uji terpisah untuk klaim akhir. \\
\bottomrule
\end{tabular}
\end{center}

\subsection{Urutan praktik yang disarankan}
\begin{enumerate}
\item \textbf{Periksa data.} Hitung cerita, panjang byte, duplikat,
dan pembagian train/validasi. Baca beberapa contoh. Pastikan validasi
berisi cerita utuh yang tidak ada di train.
\item \textbf{Jalankan smoke hanya untuk memeriksa kode.} Dua langkah
tidak cukup untuk menilai bahasa.
\item \textbf{Latih dense lebih dulu.} Simpan kurva loss, ringkasan,
dan checkpoint. Gunakan prompt dan suhu yang tetap saat membandingkan
contoh generasi.
\item \textbf{Periksa generalisasi.} Lihat apakah loss validasi masih
turun. Jangan menyimpulkan kemampuan bercerita dari loss saja; baca
apakah tokoh, peristiwa, dan akhir cerita saling berhubungan.
\item \textbf{Ubah satu variabel.} Ulangi beberapa seed jika biaya
memungkinkan. Varian sparse harus dibandingkan pada data dan anggaran
training yang sama.
\item \textbf{Simpan catatan.} Tulis hipotesis, konfigurasi, metrik,
dan contoh keluaran. Hasil yang gagal juga berguna jika penyebabnya
jelas.
\end{enumerate}

\analogi{Mengembangkan model seperti mencoba resep. Jika Anda mengganti
tepung, suhu oven, dan waktu panggang sekaligus, Anda tidak tahu
perubahan mana yang membuat hasil lebih baik atau buruk.}

\subsection{Tiga batas yang sering menipu}
\begin{enumerate}
\item \textbf{Data kecil.} Model dapat menghafal susunan suku kata dan
gaya tanpa belajar menyusun cerita baru. Bahkan 60 cerita tambahan
masih korpus yang sangat kecil untuk sekitar 1,9 juta parameter.
\item \textbf{Konteks byte pendek.} Jendela 256 byte sering hanya
memuat sebagian paragraf. Model sulit menjaga hubungan jauh antara
awal dan akhir cerita jika awalnya sudah keluar dari jendela.
\item \textbf{Metrik token bukan metrik cerita.} Loss dan perplexity
mengukur prediksi token. Cerita yang runtut perlu pemeriksaan isi:
apakah tokoh konsisten, sebab-akibat jelas, dan akhir terkait
dengan awal.
\end{enumerate}

\section{Perintah dan keluaran yang perlu dikenal}

Perintah berikut adalah contoh untuk dijalankan pengguna dari root repo.
Dokumen ini sendiri tidak menjalankan training.
\begin{lstlisting}[language={}]
# Pemeriksaan kecil
.venv/bin/python -m unittest discover -s tests
.venv/bin/python -m small_llm.compare --config configs/smoke.json --overwrite

# Satu model dense dengan korpus tambahan
.venv/bin/python -m small_llm.train \
  --config configs/stories_expanded_m1.json --variant dense

# Lihat metrik dan coba checkpoint
cat artifacts/stories_expanded_m1/dense/summary.json
cat artifacts/stories_expanded_m1/dense/metrics.jsonl
.venv/bin/python -m small_llm.generate \
  --checkpoint artifacts/stories_expanded_m1/dense/model.pt \
  --prompt "Pagi itu, Mira menemukan "
\end{lstlisting}
Folder hasil yang sudah ada tidak ditimpa tanpa \code{--overwrite}.
Opsi itu berarti mengganti hasil lama, sehingga sebaiknya simpan
folder eksperimen yang ingin dipertahankan.

\section{Ringkasan hubungan rumus dan kode}
\begin{longtable}{p{0.25\linewidth}p{0.34\linewidth}>{\raggedright\arraybackslash}p{0.31\linewidth}}
\toprule
Konsep & Rumus inti & Lokasi kode \\
\midrule
\endhead
Byte dan EOS & $\mathcal V=\{0,\ldots,256\}$ &
\code{tokenizer.py}: \code{encode}, \code{decode} \\
Batch bergeser & $Y_i=X_{i+1}$ &
\code{data.py}: \code{sample_batch} \\
Embedding & $\mathbf x_i=E_t[t_i]+E_p[i]$ &
\code{model.py}: \code{StoryTransformer.forward} \\
Attention causal & $\operatorname{softmax}(QK^\top/\sqrt{d_h}+M)V$ &
\code{model.py}: \code{CausalAttention} \\
MLP dense & $W_2\operatorname{ReLU}(W_1x+b_1)+b_2$ &
\code{model.py}: \code{DenseMLP} \\
Loss token & $-\operatorname{mean}\log p(y)$ &
\code{model.py}: \code{F.cross_entropy} \\
Mask sparse & $W_{\mathrm{efektif}}=W\odot M$ &
\code{sparse.py}: \code{PrunedMLP}, \code{SETMLP} \\
TopK neuron & $\sum_{j\in S}h_jW_{2,:,j}+b_2$ &
\code{sparse.py}: \code{TopKMLP} \\
MoE & $k\sum_{e\in S}\alpha_e E_e(x)$ &
\code{sparse.py}: \code{MoEMLP} \\
Akumulasi gradien & $G^{-1}\sum_j\nabla\mathcal L_j$ &
\code{train.py}: \code{train_variant} \\
Sampling suhu & $\operatorname{softmax}(z/T)$ &
\code{model.py}: \code{generate} \\
\bottomrule
\end{longtable}

\paragraph{Akhir panduan.}
Untuk mengembangkan model, mulailah dari peta pada Bagian
\ref{sec:peta}, ikuti satu batch melalui Bagian \ref{sec:data} dan
\ref{sec:model}, lalu gunakan Bagian \ref{sec:pengembangan} untuk
merancang satu percobaan yang dapat diperiksa kembali.

\end{document}
\section{Langkah 5: evaluasi dan berkas hasil}
\label{sec:evaluasi}

\subsection{Evaluasi lokal: \code{evaluate} dalam \code{train.py}}

Model masuk ke mode \code{eval} dan berjalan tanpa gradien. Fungsi ini
mengambil \code{eval_batches} jendela acak dari \textbf{cerita validasi},
kemudian menghitung rata-rata cross-entropy:
\[
\mathcal L_{\mathrm{valid, lokal}}
=\frac{1}{J}\sum_{j=1}^{J}\mathcal L_j,
\qquad J=\code{eval_batches}.
\]
Generator titik awal selalu dibuat lagi dengan \code{seed + 2}. Artinya,
pada setiap waktu evaluasi, jendela validasi yang dipakai sama. Perubahan
loss lebih mudah ditafsirkan karena pertanyaannya tidak berganti-ganti.
Namun yang diukur tetap \emph{sampel jendela}, bukan semua byte validasi.

\analogi{Murid mengerjakan set soal ujian yang sama setiap beberapa jam.
Kemajuannya terlihat lebih jelas daripada bila soalnya selalu diacak,
tetapi hasil itu belum mewakili semua kemungkinan soal.}

Baris langkah 0 menyimpan loss validasi \emph{sebelum} pembaruan bobot.
\code{eval_every} menentukan interval evaluasi berikutnya; langkah terakhir
selalu dievaluasi. \code{log_every} menentukan kapan baris training dicatat.
Pada baris tanpa evaluasi, \code{valid_loss} bernilai \code{null}, bukan
angka nol.

\subsection{Cara membaca \code{metrics.jsonl} dan \code{summary.json}}

\code{metrics.jsonl} berisi satu objek JSON per baris. Misalnya:
\begin{lstlisting}[language={}]
{"step": 100, "train_loss": 2.20, "valid_loss": 2.30,
 "step_seconds": 0.06}
\end{lstlisting}
\code{summary.json} meringkas nilai terakhir, jumlah parameter, pembagian
berkas, jumlah token, waktu, dan memori. \code{model.pt} berisi bobot,
konfigurasi arsitektur, konfigurasi sparsity, dan nama varian. Checkpoint
lokal \textbf{tidak} menyimpan state AdamW atau posisi generator acak;
menjalankan training lagi berarti memulai percobaan baru, bukan meneruskan
tepat dari langkah terakhir.

\begin{center}
\begin{tabular}{p{0.32\linewidth}p{0.59\linewidth}}
\toprule
Metrik & Cara membaca \\
\midrule
\code{train_loss} & Loss token pada microbatch langkah terakhir;
berfluktuasi karena batch berubah. \\
\code{valid_loss} & Loss pada jendela validasi yang dipilih dengan seed
tetap; bandingkan dalam eksperimen yang datanya sama. \\
\code{initial_valid_loss} & Titik awal sebelum training. \\
\code{train_tokens_per_second} & $BGL\times\code{steps}$ dibagi waktu
langkah training; evaluasi dan generasi tidak dihitung. \\
\code{peak_process_rss_mb} & Puncak memori proses Python. Ini bukan ukuran
lengkap semua memori GPU atau biaya produksi. \\
\bottomrule
\end{tabular}
\end{center}

Jika \code{train_loss} turun tetapi \code{valid_loss} naik terus, model
mungkin terlalu menyesuaikan diri pada data train. Jika keduanya turun,
model belajar memprediksi token pada distribusi yang diuji. Kualitas
cerita masih harus diperiksa dengan membaca contoh baru.

Perplexity per token dapat dihitung dari cross-entropy:
\[
\operatorname{PPL}=\exp(\mathcal L).
\]
Karena token di sini adalah \textbf{byte}, ini adalah perplexity per byte.
Jangan bandingkan langsung dengan angka dari tokenizer kata atau subkata.
Program lokal tidak menyimpan perplexity sebagai kolom terpisah.

\subsection{Perbandingan: \code{compare.py}}

\code{compare.py} menjalankan dense dan empat varian sparse dengan
konfigurasi yang sama. Secara bawaan tiap varian memakai proses Python
terpisah, agar puncak RSS satu varian tidak terbawa ke pengukuran varian
lain. Setelah itu program mengumpulkan \code{summary.json} menjadi
\code{comparison.csv}. Opsi \code{--variants} memilih sebagian varian;
\code{--overwrite} diperlukan jika folder output varian sudah ada dan
pengguna memang ingin menggantinya. Dengan \code{--in-process}, varian
dijalankan dalam satu proses dan angka puncak RSS per varian dikosongkan.

\perhatikan{Perbandingan adil memerlukan data, pembagian train/validasi,
seed, batch token, jumlah langkah, optimizer, dan perangkat yang sama.
Jumlah parameter serta operasi tiap varian tetap berbeda. Loss dari
\code{stories_demo} tidak langsung sebanding dengan loss dari
\code{stories_expanded} atau dataset Frog karena data ujinya berubah.}

\section{Langkah 6: model menulis lanjutan}
\label{sec:generasi}

\code{generate.py} memuat \code{model.pt} dengan
\code{torch.load(..., weights_only=True)}, membangun kembali Transformer
dan jenis MLP dari metadata, lalu memuat \code{state_dict}. Prompt diubah
menjadi byte, kemudian \code{StoryTransformer.generate} mengulang:
\begin{enumerate}
\item Ambil paling banyak $L$ token terakhir sebagai konteks.
\item Hitung logit pada posisi terakhir.
\item Bagi logit dengan suhu $T>0$, ubah menjadi probabilitas,
pilih satu token, lalu tambahkan ke konteks.
\item Berhenti jika ID 256 terpilih atau batas token baru tercapai.
\end{enumerate}
Distribusi sampling pada suhu $T$ adalah
\[
p_T(v)=
\frac{\exp(z_v/T)}
{\sum_{u\in A}\exp(z_u/T)},\qquad v\in A,
\]
dengan $A$ himpunan token yang diizinkan. Suhu lebih rendah membuat
pilihan lebih terkonsentrasi pada logit tinggi; suhu lebih tinggi
membuat hasil lebih bervariasi. Mengubah suhu \textbf{tidak} mengajari
model fakta atau alur baru.

\subsection{Mengapa keluaran tetap UTF-8 tetapi bisa aneh?}

Karena satu karakter Unicode dapat memakai beberapa byte, sampling byte
secara bebas bisa menghasilkan urutan yang rusak. Implementasi lokal
melacak apakah model sedang berada di tengah karakter UTF-8. Jika ya,
hanya byte lanjutan yang sah diizinkan. Di akhir batas token, model
tidak boleh memulai karakter yang tidak akan sempat diselesaikan.
Karakter pengganti akibat byte rusak pun dapat dihindari.

\analogi{Untuk menulis satu huruf yang terdiri dari tiga potong puzzle,
model harus menyelesaikan potong kedua dan ketiga sebelum memulai huruf
baru. Aturan ini hanya menjaga bentuk tulisannya, bukan isi ceritanya.}

Generasi menampilkan kembali prompt diikuti kelanjutannya. Batas
\code{--max-new-tokens} menghitung \textbf{byte baru}, bukan kata atau
kalimat. Model 300 langkah pada 12 cerita dapat membuat suku kata yang
menyerupai bahasa Indonesia tetapi belum menghasilkan cerita runtut.

\section{Tes: apa yang dibuktikan \code{test_small_llm.py}?}
\label{sec:tes}

\begin{enumerate}
\item \textbf{Tokenizer dan data:} teks UTF-8 dapat dikodekan dan dibaca
kembali; berkas train dan validasi tidak tumpang tindih; target adalah
input yang bergeser satu token.
\item \textbf{Causal attention:} mengubah token masa depan tidak mengubah
logit posisi sebelumnya.
\item \textbf{TopK dan MoE:} dalam pengaturan yang memakai semua
neuron/expert, keluaran dapat disamakan dengan MLP dense sesuai
konstruksi tes.
\item \textbf{SET:} rewiring menjaga jumlah koneksi aktif dan menghapus
state optimizer pada posisi yang berubah.
\item \textbf{Checkpoint:} arsitektur sparse dapat dibangun kembali
dan menghasilkan keluaran yang sama sesudah bobot dimuat.
\item \textbf{Training singkat:} setiap varian dapat menjalankan satu
langkah dan menyimpan checkpoint dengan loss hingga.
\item \textbf{Generasi:} byte baru mengikuti urutan UTF-8 yang valid
pada contoh yang sengaja memberi skor tinggi pada byte tidak sah.
\end{enumerate}

\perhatikan{Tes ini membuktikan sifat kode pada contoh kecil. Tes satu
langkah tidak membuktikan bahwa model mampu bercerita. Kualitas bahasa
memerlukan data dan evaluasi yang terpisah.}


